\clearpage
\VerdierPhysicalPageAnchor{0143}%
\begingroup
\small
\setlength{\parskip}{0pt}%
\setlength{\emergencystretch}{1.5em}%

\begingroup\itshape
\VerdierSemanticAnchor{ch2-2-3-1-continuation}%
对于每个满足
$\mathcal B\subset\mathcal A\subset\mathcal D$ 的满三角子范畴
$\mathcal A$，记 $S_{\mathcal D}(\mathcal B)$（分别地，
$S_{\mathcal A}(\mathcal B)$）为 $\mathcal D$（分别地，$\mathcal A$）中
与 $\mathcal B$ 关联的乘法系 (2.1.7)。

\VerdierSemanticAnchor{ch2-2-3-1-a-intersection}
\par\noindent a) 有
$S_{\mathcal D}(\mathcal B)\cap\operatorname{Fl}(\mathcal A)
=S_{\mathcal A}(\mathcal B)$。

\VerdierSemanticAnchor{ch2-2-3-1-b-fully-faithful}
\par\noindent b) 典范函子
$\mathcal A/\mathcal B\longrightarrow\mathcal D/\mathcal B$ 充分忠实，
并且在对象上单射。这个函子的像是 $Q(\mathcal A)$，其中
$Q\colon\mathcal D\longrightarrow\mathcal D/\mathcal B$ 是典范函子。

\VerdierSemanticAnchor{ch2-2-3-1-c-quotient-transitivity}
\par\noindent c) 典范函子
$\mathcal D/\mathcal A\longrightarrow
(\mathcal D/\mathcal B)/(\mathcal A/\mathcal B)$ 是范畴同构。

\VerdierSemanticAnchor{ch2-2-3-1-d-saturation}
\par\noindent d) 若范畴 $\mathcal A$ 是饱和的 (2.1.6)，则子范畴
$Q(\mathcal A)$ 是饱和的。

\VerdierSemanticAnchor{ch2-2-3-1-inverse-image-setup}
对于 $\mathcal D/\mathcal B$ 的每个满三角子范畴 $\mathcal C$，记
$Q^{-1}(\mathcal C)$ 为 $\mathcal D$ 的如下满子范畴：其对象是
$\mathcal D$ 中满足 $Q(X)$ 同构于范畴 $\mathcal C$ 的某个对象的对象
$X$。

\VerdierSemanticAnchor{ch2-2-3-1-a-bis}
\par\noindent a)$^{\mathrm{bis}}$ 范畴 $Q^{-1}(\mathcal C)$ 是
$\mathcal D$ 的严格满三角子范畴。

\VerdierSemanticAnchor{ch2-2-3-1-b-bis}
\par\noindent b)$^{\mathrm{bis}}$ 乘法系
$S_{\mathcal D}(Q^{-1}(\mathcal C))$ 由 $\mathcal D$ 中经 $Q$ 的像属于
$S_{\mathcal D/\mathcal B}(\mathcal C)$ 的态射组成。

\VerdierSemanticAnchor{ch2-2-3-1-c-bis}
\par\noindent c)$^{\mathrm{bis}}$ 子范畴 $Q^{-1}(\mathcal C)$ 是饱和的，
当且仅当子范畴 $\mathcal C$ 是饱和的。

\VerdierSemanticAnchor{ch2-2-3-1-d-bis}
\par\noindent d)$^{\mathrm{bis}}$ 映射
$\mathcal A\longmapsto Q(\mathcal A)$ 给出一个有序集同构：它把按包含关系
排序的、包含 $\widehat{\mathcal B}$（函子 $Q$ 的核）的 $\mathcal D$ 的严格
满三角子范畴所成的集合，映到按包含关系排序的 $\mathcal D/\mathcal B$ 的
严格满三角子范畴所成的集合。逆映射为
$\mathcal C\longmapsto Q^{-1}(\mathcal C)$。
\endgroup

\VerdierSemanticAnchor{ch2-2-3-1-proof-ab}
第一个断言是显然的。第二个断言立即由范畴
$\mathcal A/\mathcal B$ 与 $\mathcal D/\mathcal B$ 中态射的描述 (2.2.1)
得到。

\VerdierSemanticAnchor{ch2-2-3-1-proof-c}
为证明 c)，只需引用商范畴的泛性质 (2.2.11)。

\VerdierSemanticAnchor{ch2-2-3-1-proof-d}
来证明 d)。把 $\mathcal A$ 替换为它的严格满闭包只会在等价意义下改变
$Q(\mathcal A)$，因此可假设 $\mathcal A$ 严格满。这时函子
$\mathcal D\longrightarrow\mathcal D/\mathcal A$ 的核是范畴 $\mathcal A$
(2.2.11)。该函子分解为
$\mathcal D\longrightarrow\mathcal D/\mathcal B
\longrightarrow\mathcal D/\mathcal A$。由此可知，典范函子
$\mathcal D/\mathcal B\longrightarrow\mathcal D/\mathcal A$ 的核是
$Q(\mathcal A)$。因而 $Q(\mathcal A)$ 是饱和的 (2.1.16)。

\VerdierSemanticAnchor{ch2-2-3-1-proof-a-bis-b-bis}
断言 a)$^{\mathrm{bis}}$ 是显然的，断言 b)$^{\mathrm{bis}}$ 也是如此。

\VerdierSemanticAnchor{ch2-2-3-1-proof-c-bis-forward}
假设 $\mathcal C$ 是饱和的。还可假设 $\mathcal C$ 严格满，因为若以
$\widehat{\mathcal C}$ 表示 $\mathcal C$ 的严格满闭包，则范畴
$\widehat{\mathcal C}$ 严格满且饱和，并且
$Q^{-1}(\widehat{\mathcal C})=Q^{-1}(\mathcal C)$。于是范畴
$\mathcal C$ 是某个正合函子的核 (2.2.11)。因此，
$Q^{-1}(\mathcal C)$ 是某个正合函子的核，从而是饱和的 (2.1.16)。

\VerdierSemanticAnchor{ch2-2-3-1-p144-boundary}
假设 $Q^{-1}(\mathcal C)$ 是
\endgroup
